\section{Related Work}
\label{sec:related}

\paragraph{Keystroke Dynamics.}
Keystroke dynamics as a behavioral biometric were first studied for
authentication in the 1980s~\cite{monrose2000keystroke}.
Killourhy and Maxion~\cite{killourhy2009comparing} established a benchmark
dataset and systematic comparison of 14 classifiers, finding that simple
distance-based methods often match more complex approaches.
Subsequent work introduced deep learning for keystroke authentication,
including RNNs for sequence modeling and Siamese networks for one-shot
verification~\cite{lu2019deep}.
All of these approaches optimize for \emph{cross-user} discrimination.
Our work targets the orthogonal goal of \emph{within-user} longitudinal drift
detection, for which population-template methods are fundamentally ill-suited.

\paragraph{Anomaly Detection with Autoencoders.}
Autoencoders were introduced for nonlinear dimensionality reduction by
Rumelhart \etal~\cite{rumelhart1986learning} and popularized for deep
representations by Hinton and Salakhutdinov~\cite{hinton2006reducing}.
Their application to anomaly detection exploits the observation that a model
trained exclusively on normal data will reconstruct normal inputs well but
yield high error on anomalous inputs~\cite{an2015variational}.
Vincent \etal~\cite{vincent2010stacked} introduced denoising autoencoders,
training to reconstruct clean inputs from corrupted versions, which improves
robustness and forces learning of more meaningful representations.
We adopt a denoising variant with reproducible 10\% input masking.

\paragraph{Differential Privacy in Deep Learning.}
Dwork \etal~\cite{dwork2006calibrating} formalized differential privacy (DP)
as a mathematically rigorous definition of privacy for statistical algorithms.
Abadi \etal~\cite{abadi2016deep} extended DP to neural network training via
DP-SGD: gradient clipping followed by calibrated Gaussian noise addition,
with privacy budget tracked via moments accountant.
The Opacus library~\cite{yousefpour2021opacus} provides an efficient
PyTorch implementation with R\'{e}nyi Differential Privacy (RDP) accounting.
In the context of personal behavioral data — where lessons encode daily
behavioral patterns — formal privacy guarantees are especially important.

\paragraph{Longitudinal Behavioral Monitoring.}
Digital biomarker research has demonstrated that passive sensing of
device-interaction patterns can detect cognitive fluctuations~\cite{dagum2018digital}.
Typing-based monitoring has been proposed for Parkinson's disease~\cite{equi2021detection}
and attention deficit detection~\cite{arroyo2009user}.
These systems typically rely on engineered statistical features rather than
learned representations.
Our work contributes a learned, privacy-preserving representation that
adapts entirely to the individual without supervision.

\paragraph{N-of-1 Designs.}
N-of-1 randomized controlled trials~\cite{lillie2011nof1} treat each participant
as their own control across multiple crossover periods, enabling causal inference
at the individual level.
The philosophy has been applied to personalized medicine and digital therapeutics.
Our work is, to our knowledge, the first application of the N-of-1 design
philosophy to a deep autoencoder for keystroke behavioral monitoring.
